
\section{Background}

\subsection{Singular Value Decomposition}

For any matrix $W \in \mathbb{R}^{m \times n}$, the SVD is:

\[
W = U \Sigma V^\top
\tag{3.1}
\]

where $U \in \mathbb{R}^{m \times m}$, $V \in \mathbb{R}^{n \times n}$ are 
orthogonal, and $\Sigma = \mathrm{diag}(\sigma_1, \dots, \sigma_{\min(m,n)})$ 
with $\sigma_1 \geq \sigma_2 \geq \dots \geq 0$.

\textbf{Truncated SVD.} The best rank-$r$ approximation in Frobenius norm is:

\[
W_r = U_{:, :r} \Sigma_{:r, :r} V_{:r, :}^\top
\tag{3.2}
\]

with $\|W - W_r\|_F^2 = \sum_{i > r} \sigma_i^2$.

\textbf{Symmetric factorization.} Distributing $\Sigma^{1/2}$ between factors:

\[
U_r = U_{:, :r} \Sigma_{:r, :r}^{1/2}, \quad
V_r = \Sigma_{:r, :r}^{1/2} V_{:r, :}^\top
\tag{3.3}
\]

ensures $\|U_r\|_2 = \|V_r\|_2 = \sigma_1^{1/2}$, giving balanced magnitudes 
and improved numerical stability for FP16 computation.

\subsection{Cumulative Squared Energy}

Define the cumulative squared energy at rank $r$:

\[
\mathcal{E}(r) = \frac{\sum_{i=1}^{r} \sigma_i^2}{\sum_{j=1}^{\min(m,n)} \sigma_j^2}
\tag{3.4}
\]

Properties:
\begin{itemize}
\item $\mathcal{E}(0) = 0$
\item $\mathcal{E}$ is monotone increasing
\item $\mathcal{E}(\min(m, n)) = 1$
\item $\mathcal{E}(r)$ measures the fraction of $\|W\|_F^2$ captured by the top-$r$ directions

\end{itemize}
\textbf{Energy-thresholded rank.} Given $\tau \in (0, 1)$:

\[
r^*(\tau) = \min\{r : \mathcal{E}(r) \geq \tau\}
\tag{3.5}
\]

At this rank:

\[
\|W - W_{r^*}\|_F^2 = (1 - \mathcal{E}(r^*)) \|W\|_F^2 \leq (1 - \tau) \|W\|_F^2
\tag{3.6}
\]

so relative Frobenius error is $\leq \sqrt{1 - \tau}$.

\subsection{Spectral Decay in Attention Weights}

Attention projection matrices exhibit heterogeneous spectral decay:

\begin{itemize}
\item \textbf{Query projections:} Moderate decay. Effective rank ~40--80\% of full.
\item \textbf{Key projections:} Slower decay (keys spread information across dimensions).
\item \textbf{Value projections:} Faster decay (values have lower effective rank).
\item \textbf{Output projections:} Moderate decay; depends on head-combining structure.

\end{itemize}
\textbf{Figure 3.1} shows singular value spectra for DistilBERT layer 0 vs layer 3. 
The difference in decay rate justifies per-layer rank selection.

\subsection{LoRA-Style Residual Adapters}

LoRA (Hu et al., 2021) parameterizes a weight update as:

\[
\Delta W = \frac{\alpha}{r_L} B A
\tag{3.7}
\]

with $A \in \mathbb{R}^{r_L \times d_\text{in}}$, $B \in \mathbb{R}^{d_\text{out} \times r_L}$. 
The initialization $B = 0$ ensures $\Delta W = 0$ at step 0.

\textbf{Applied to factorized layers.} The full forward:

\[
y = x (U V)^\top + \frac{\alpha}{r_L} x (B A)^\top
\tag{3.8}
\]

At initialization, $y = x (U V)^\top$, matching the frozen base.

\subsection{T5 Relative Attention Bias}

T5 (Raffel et al., 2020) adds a learnable bias to attention scores based on 
the bucketed relative position of query and key:

\[
b(i, j) = \text{Embed}[\text{bucket}(j - i)]
\tag{3.9}
\]

The bucket function is:

\[
\text{bucket}(d) = \begin{cases}
d & \text{if } |d| < \text{max\_exact} \\
\text{max\_exact} + \text{sign}(d) \cdot \lfloor \log(|d| / \text{max\_exact}) \cdot \text{scale} \rfloor & \text{otherwise}
\end{cases}
\tag{3.10}
\]

with $\text{max\_exact} = 8$, $\text{num\_buckets} = 32$. The bias is added 
to attention scores before softmax. Adaptive SVD preserves this bias.

\subsection{Grouped-Query Attention (GQA)}

Modern models (Qwen, LLaMA) use GQA where $H_{kv} < H$:

\[
K, V \in \mathbb{R}^{B \times L \times H_{kv} d_h}, \quad H_\text{repeat} = H / H_{kv}
\tag{3.11}
\]

The KV heads are repeated via \texttt{repeat\_kv} before attention. Adaptive SVD 
preserves GQA structure.

\subsection{Rotary Position Embeddings (RoPE)}

RoPE (Su et al., 2021) rotates $Q, K$ by position-dependent angles:

\[
\tilde{Q}_i = Q_i \cos(\theta_i) + \mathrm{rot}(Q_i) \sin(\theta_i)
\tag{3.12}
\]

Adaptive SVD applies RoPE after the factorized $Q, K$ projections.

\subsection{Notation Summary}
